\subsection{Closing Comment}\label{sec:ClosingComment}
\keeplisting{5}
A closing comment like this, for a Java source file,
\begin{lstlisting}[caption={Closing Comment for a Java file}]
/*
 *  End of File
 */
\end{lstlisting}
at the end of the source file is helpful to detect whether a source file was corrupted: if missing, there is a good chance that the file is incomplete.

I confess that this kind of error has got very unlikely with modern hard-drives and SSDs, but it can still happen when sending around sources over the net.

Below samples of closing comments for some of the most common file types that appear as sources within a Java project, in addition to \verb#*.java# files; just be creative for other file types.

\keeplisting{5}
For a resource file\index{resource file}\index{properties file}, the closing comment has the following form: 
\begin{lstlisting}[language=,caption={Closing Comment for a resource or a script file}]
#
# End of File
#
\end{lstlisting}

\keeplisting{5}
And for an XML\index{XML} or HTML\index{HTML} file:
\begin{lstlisting}[language=XML,caption={Closing Comment for a HTML or an XML file}]
<!--
End of File
-->
\end{lstlisting}

\subsubsection{Principles}
\begin{enumerate}[label=P\arabic*.]
\item{Each source file must have a closing comment, where feasible.}
\end{enumerate}

%==============================================================================
% End of file!!
%==============================================================================
